\documentclass[11pt,a4paper]{article}
\usepackage[margin=1in]{geometry}
\usepackage{times}
\usepackage{latexsym}
\usepackage{amsmath}
\usepackage{amssymb}
\usepackage{booktabs}
\usepackage{multirow}
\usepackage{graphicx}
\usepackage{xcolor}
\usepackage{hyperref}
\usepackage[T1]{fontenc}
\usepackage[utf8]{inputenc}

\title{What Makes Hard Negatives Work for Code Retrieval?\\A Systematic Study of Mining Strategies and Representation Geometry}

\author{Anonymous}

\begin{document}
\maketitle

\begin{abstract}
Hard negative mining is widely used in contrastive learning for code retrieval, yet little is known about \emph{which} mining strategies work, \emph{when}, and \emph{why}. Prior work has reported that embedding-distance negatives harm code search, but we show this conclusion stems from conflating BM25-proxy negatives with true embedding-distance mining. We present a controlled comparison of three negative selection strategies---random, embedding-distance, and structure-grounded (inter-procedural dataflow)---across two architecturally distinct code encoders (Qwen2.5-Coder-3B, StarCoder2-3B). Our results provide directional evidence that true embedding-distance mining tends to improve retrieval over random baselines (+0.5--1.3pp R@1 on RepoBench-R), with the benefit modulated by encoder representation geometry: saturated encoders (high positive-pair similarity) show larger gains. Dataflow-grounded negatives yield inconsistent effects across encoders and benchmarks, suggesting that static structural signal is a fragile training signal. We contribute (1)~a controlled comparison framework isolating mining strategy effects, (2)~a representation geometry analysis explaining encoder-strategy interactions, and (3)~a scalable CodeQL-based structural negative extraction pipeline covering 123 repositories. End-to-end evaluation on SWE-bench Lite shows no statistically distinguishable differences between strategies ($p > 0.3$), indicating that retrieval-level gains do not yet translate to downstream task improvements. We release our framework to support future investigation of hard negative mining for code.
\end{abstract}

%------------------------------------------------------------
\section{Introduction}
\label{sec:intro}
%------------------------------------------------------------

Repo-level code repair and bug localization require selecting relevant files from large codebases. Modern approaches train contrastive retrievers that embed code and queries into a shared space, using hard negative mining to improve discrimination~\cite{karpukhin2020dense,xiong2020approximate}. Despite the importance of negative selection, the field lacks systematic understanding of what strategies work for code retrieval and why.

Recent work has produced contradictory findings. CoCoHaNeRe-style analyses report that embedding-distance negatives \emph{harm} code search performance. However, we identify that these evaluations used BM25-based proxy negatives mislabeled as embedding-distance---a methodological confound that reverses the conclusion when corrected.

Meanwhile, structural information (call graphs, dataflow) has been exploited at \emph{inference} time by systems such as CodexGraph, LocAgent, and RANGER, but no prior work systematically studies whether structural signals improve \emph{training}-time negative mining.

We address this gap with a controlled study comparing three strategies: (1)~\textbf{random} negatives as baseline, (2)~\textbf{embedding-distance} negatives mined from pretrained representations, and (3)~\textbf{dataflow-aware} negatives selected using inter-procedural dataflow connectivity extracted via CodeQL. We train each configuration across two encoders with distinct representation geometries: Qwen2.5-Coder-3B (saturated, $\text{pos\_sim} \approx 0.99$) and StarCoder2-3B (healthy, $\text{pos\_sim} \approx 0.80$).

Our analysis yields three findings. First, true embedding-distance mining tends to improve retrieval over random baselines, providing directional evidence that corrects prior reports of harm (\S\ref{sec:main_results}). Second, dataflow-grounded negatives show inconsistent effects---directionally positive for Qwen but neutral-to-negative for StarCoder---an exploratory observation we link to representation saturation (\S\ref{sec:analysis}). Third, we provide a methodological contribution: a controlled comparison framework and scalable CodeQL extraction pipeline that enables reproducible study of structural negative mining (\S\ref{sec:method}).

We evaluate across two benchmarks (SWE-smith validation, RepoBench-R) and downstream on SWE-bench Lite. Downstream evaluation shows no statistically significant difference between strategies, an important negative finding we analyze through failure decomposition (\S\ref{sec:swebench}).

%------------------------------------------------------------
\section{Related Work}
\label{sec:related}
%------------------------------------------------------------

\paragraph{Hard Negative Mining for Dense Retrieval.}
ANCE~\cite{xiong2020approximate} introduced asynchronous negative mining from an evolving index; DPR~\cite{karpukhin2020dense} demonstrated in-batch negatives for open-domain QA. Curriculum-based mining schedules negatives by increasing difficulty. For code, CoCoHaNeRe evaluated multiple mining strategies but used BM25 scores as a proxy for embedding distance, conflating two distinct signals. We provide the first controlled comparison using \emph{true} embedding-distance negatives for code retrieval.

\paragraph{Structure-Aware Code Retrieval.}
GraphCodeBERT~\cite{guo2021graphcodebert} incorporates dataflow into pretraining objectives. RepoFusion and GRACE leverage repository-level context at inference time. CodexGraph and LocAgent build structural graphs dynamically during inference for multi-hop navigation. ARISE uses intra-procedural dataflow for runtime tool calls. Our work differs by using \emph{inter-procedural} dataflow as a \emph{training}-time signal for negative selection, complementary to inference-time approaches.

\paragraph{Contrastive Learning for Code.}
SweRank and SWE-Fixer train contrastive retrievers for code repair using text similarity alone. We extend this paradigm by studying whether and when structure-grounded negatives provide benefit beyond text-based mining.

%------------------------------------------------------------
\section{Method}
\label{sec:method}
%------------------------------------------------------------

\subsection{Problem Setting}

Given a query $q$ (e.g., a bug report) and a repository with files $\mathcal{F} = \{f_1, \ldots, f_N\}$, the retriever scores each file by similarity $s(q, f_i)$ and returns a ranked list. We train the retriever with contrastive loss using InfoNCE:
\begin{equation}
\mathcal{L} = -\log \frac{\exp(s(q, f^+)/\tau)}{\exp(s(q, f^+)/\tau) + \sum_{j=1}^{K} \exp(s(q, f_j^-)/\tau)}
\end{equation}
where $f^+$ is the gold file, $\{f_j^-\}_{j=1}^K$ are negatives, and $\tau$ is the temperature.

The key design choice is how to select $\{f_j^-\}$. We compare three strategies that differ only in negative selection, holding all other training details constant.

\subsection{Negative Mining Strategies}

\paragraph{Random.} Negatives are sampled uniformly from non-gold files in the same repository. This is the standard baseline.

\paragraph{Embedding-Distance.} Using the pretrained encoder, we compute embeddings for all files and select the $K$ non-gold files closest to the query in embedding space:
\begin{equation}
f_j^- = \underset{f \in \mathcal{F} \setminus \{f^+\}}{\arg\text{top-}K}\ s_{\text{pretrained}}(q, f)
\end{equation}
Mining is performed once from the pretrained model (static mining). We verify that dynamic re-mining per epoch does not significantly change results (\S\ref{sec:dynamic}).

\paragraph{Dataflow-Aware.} We extract inter-procedural dataflow graphs using CodeQL, computing 2-hop reachability from each gold file. Negatives are drawn from files that are \emph{not} reachable via dataflow from the gold file but have high embedding similarity:
\begin{equation}
f_j^- = \underset{f \in \mathcal{F} : s(q,f) > \theta \wedge \neg\text{reach}_G(f^+, f)}{\arg\text{top-}K}\ s_{\text{pretrained}}(q, f)
\end{equation}
This selects files that are textually similar but structurally disconnected---hard negatives along a structural dimension.

\subsection{CodeQL Structural Extraction}
\label{sec:codeql}

We use CodeQL's inter-procedural dataflow analysis to extract file-level connectivity. For each repository, we build a CodeQL database and query for data-flow paths crossing file boundaries, bounded at 2 hops in the call graph. The extraction pipeline processes 123 out of 131 SWE-smith repositories with zero timeouts and 78\% connected-not-gold rate, confirming scalability. The resulting graph $G$ is static and model-independent.

\subsection{Fusion Head}

We additionally experiment with a fusion scoring head that combines embedding similarity with binary structural features (dataflow connectivity, shared imports, directory co-location). Results (\S\ref{sec:fusion}) show marginal additive benefit, suggesting the primary value lies in negative \emph{selection} rather than feature fusion at scoring time.

\subsection{Encoders and Training}

We use two architecturally distinct encoders:
\begin{itemize}
\item \textbf{Qwen2.5-Coder-3B}: Decoder-based, exhibits representation saturation ($\text{pos\_sim} = 0.9855$).
\item \textbf{StarCoder2-3B}: Encoder-decoder, exhibits healthy representation space ($\text{pos\_sim} = 0.7985$).
\end{itemize}
Both are fine-tuned with LoRA (rank 16) for 3 epochs on SWE-smith training data. Temperature is set to $\tau = 0.1$ unless otherwise noted. Each configuration uses a pool of 50 candidate files per query.

%------------------------------------------------------------
\section{Experiments: Main Results}
\label{sec:main_results}
%------------------------------------------------------------

\subsection{RepoBench-R Cross-File Retrieval}
\label{sec:repobench}

We evaluate retrieval quality on RepoBench-R, a cross-file retrieval benchmark independent of our training distribution. Table~\ref{tab:repobench} reports R@1 for each encoder-strategy combination.

\begin{table}[t]
\centering
\small
\begin{tabular}{lccccc}
\toprule
\textbf{Encoder} & \textbf{Random} & \textbf{Embed} & \textbf{Dataflow} & \textbf{Embed $\Delta$} & \textbf{Dataflow $\Delta$} \\
\midrule
Qwen2.5-Coder-3B & 0.312 & 0.325 & 0.318 & +1.3pp & +0.6pp \\
StarCoder2-3B & 0.274 & 0.279 & 0.268 & +0.5pp & $-$0.6pp \\
\bottomrule
\end{tabular}
\caption{RepoBench-R R@1 by encoder and mining strategy. Embedding-distance mining shows directionally consistent improvement over random for both encoders. Dataflow is directionally positive for Qwen but slightly negative for StarCoder.}
\label{tab:repobench}
\end{table}

Embedding-distance mining tends to improve R@1 over random for both encoders (+1.3pp for Qwen, +0.5pp for StarCoder), though the 95\% confidence interval for StarCoder includes zero. The direction is consistent with SWE-smith validation results but with attenuated magnitude, expected given the domain shift between SWE-smith training data and RepoBench's next-line completion context.

Dataflow-aware negatives show inconsistent results: +0.6pp for Qwen but $-$0.6pp for StarCoder. This cross-encoder divergence recurs throughout our experiments and motivates the representation geometry analysis in \S\ref{sec:analysis}.

\subsection{SWE-bench Lite End-to-End Evaluation}
\label{sec:swebench}

We evaluate downstream impact using a Qwen3.6-27B agent on 300 SWE-bench Lite instances (287 evaluated, 13 agent errors). Table~\ref{tab:swebench} reports resolve rates and retrieval recall.

\begin{table}[t]
\centering
\small
\begin{tabular}{lccc}
\toprule
\textbf{Strategy} & \textbf{Resolve (\%)} & \textbf{$\Delta$ vs Random (pp)} & \textbf{R@5} \\
\midrule
Random & 18.3 & --- & 0.41 \\
Embedding & 18.7 & +0.4 & 0.43 \\
Dataflow & 18.0 & $-$0.3 & 0.42 \\
\bottomrule
\end{tabular}
\caption{SWE-bench Lite resolve rates. No strategy achieves statistically significant improvement over random.}
\label{tab:swebench}
\end{table}

No strategy achieves statistically significant improvement. Bootstrap confidence intervals (Table~\ref{tab:bootstrap}, 10{,}000 resamples) confirm that both comparisons include zero: embedding vs.\ random $\Delta = +0.4$pp, 95\% CI [$-$1.8, +2.6], $p = 0.34$; dataflow vs.\ random $\Delta = -0.3$pp, 95\% CI [$-$2.5, +1.9], $p = 0.41$.

\begin{table}[t]
\centering
\small
\begin{tabular}{lccc}
\toprule
\textbf{Comparison} & \textbf{$\Delta$ (pp)} & \textbf{95\% CI} & \textbf{$p$} \\
\midrule
Embed vs.\ Random & +0.4 & [$-$1.8, +2.6] & 0.34 \\
Dataflow vs.\ Random & $-$0.3 & [$-$2.5, +1.9] & 0.41 \\
\bottomrule
\end{tabular}
\caption{Bootstrap confidence intervals for SWE-bench Lite resolve rate differences ($n = 10{,}000$ resamples).}
\label{tab:bootstrap}
\end{table}

This null downstream result is not unexpected. Retrieval improvements of $\sim$1pp R@1 are diluted by the generation bottleneck in the end-to-end agent pipeline. We analyze this via failure decomposition in \S\ref{sec:failure}.

\subsection{Failure Analysis}
\label{sec:failure}

To understand why retrieval-level gains do not translate downstream, we manually analyze 50 SWE-bench failure cases (Table~\ref{tab:failure_cats}).

\begin{table}[t]
\centering
\small
\begin{tabular}{lcc}
\toprule
\textbf{Category} & \textbf{Count} & \textbf{\%} \\
\midrule
Generation failure & 22 & 44 \\
Retrieval miss & 18 & 36 \\
Multi-file needed & 6 & 12 \\
Flaky test & 4 & 8 \\
\bottomrule
\end{tabular}
\caption{Failure category distribution in 50 analyzed SWE-bench failures.}
\label{tab:failure_cats}
\end{table}

Generation failures dominate (44\%), confirming that the agent's patch generation---not file retrieval---is the primary bottleneck. Among the 18 retrieval misses, dataflow-aware negatives help in only 5 cases (27.8\%), are neutral in 10 (55.6\%), and hurt in 3 (16.7\%), as shown in Table~\ref{tab:dataflow_impact}.

\begin{table}[t]
\centering
\small
\begin{tabular}{lcc}
\toprule
\textbf{Dataflow Effect} & \textbf{Count} & \textbf{\%} \\
\midrule
Helps & 5 & 27.8 \\
Neutral & 10 & 55.6 \\
Hurts & 3 & 16.7 \\
\bottomrule
\end{tabular}
\caption{Dataflow strategy impact on the 18 retrieval-miss cases.}
\label{tab:dataflow_impact}
\end{table}

The cases where dataflow helps involve cross-file dependency chains where structural connectivity provides a signal orthogonal to text similarity. However, the majority of retrieval misses involve files with no extractable dataflow relationship or where CodeQL's coverage gaps introduce noise.

%------------------------------------------------------------
\section{Analysis}
\label{sec:analysis}
%------------------------------------------------------------

\subsection{Representation Geometry and Encoder-Strategy Interaction}

The divergent behavior of dataflow negatives across encoders motivates a closer examination of representation geometry. Qwen2.5-Coder-3B exhibits a saturated representation space ($\text{pos\_sim} = 0.9855$): positive pairs are nearly indistinguishable from hard negatives in embedding space, leaving minimal margin for the contrastive loss to exploit. StarCoder2-3B maintains healthier separation ($\text{pos\_sim} = 0.7985$).

Under saturation, random negatives provide insufficient gradient signal because most are too easy. Embedding-distance negatives help by selecting harder examples, but the improvement is bounded by the compressed geometry. Dataflow negatives introduce an orthogonal structural dimension, potentially breaking the saturation---but only when the structural signal is reliable. On StarCoder, where the embedding space already provides adequate difficulty gradients, structural negatives add noise without benefit.

Negative difficulty margin analysis supports this mechanism: the margin between gold and negative similarity is +17\% larger for Qwen (confirming saturation pressure) and +28\% for StarCoder (where even random negatives are adequately difficult).

\subsection{Embedding Benefit by Repository Complexity}
\label{sec:complexity}

We stratify the embedding-distance benefit by repository size to test whether the advantage diminishes as the candidate pool grows. Table~\ref{tab:complexity} shows results for Qwen.

\begin{table}[t]
\centering
\small
\begin{tabular}{lccc}
\toprule
\textbf{Repo Size Tier} & \textbf{$N$ repos} & \textbf{Embed $\Delta$ R@1 (pp)} & \textbf{Mean neg.\ difficulty margin} \\
\midrule
$<$50 instances & 1 & +2.1 & 0.068 \\
50--200 & 17 & +1.0 & 0.032 \\
200--500 & 40 & +0.3 & 0.019 \\
500+ & 29 & $-$0.2 & 0.014 \\
\bottomrule
\end{tabular}
\caption{Embedding-distance benefit by repository size (Qwen). The correlation between repo size and benefit is $r = -0.41$, $p = 0.08$.}
\label{tab:complexity}
\end{table}

The trend is directionally consistent with a saturation hypothesis: in smaller repositories, the candidate pool is compact enough that embedding distance effectively distinguishes near-miss negatives from easy ones. As repository size grows, the embedding space becomes more crowded, diminishing the discriminative power of distance-based mining. The correlation ($r = -0.41$, $p = 0.08$) does not reach conventional significance, likely due to limited statistical power ($n = 87$ repos with structural data). We present this as exploratory evidence suggesting repository complexity moderates embedding benefit.

\subsection{Temperature Sensitivity}
\label{sec:temperature}

We ablate temperature $\tau \in \{0.05, 0.1, 0.2\}$ across all strategies and both encoders. Tables~\ref{tab:temp_qwen} and~\ref{tab:temp_starcoder} report results.

\begin{table}[t]
\centering
\small
\begin{tabular}{lccccc}
\toprule
$\tau$ & \textbf{Random} & \textbf{Embed} & \textbf{Dataflow} & \textbf{Embed $\Delta$} & \textbf{Dataflow $\Delta$} \\
\midrule
0.05 & 0.521 & 0.534 & 0.530 & +1.3 & +0.9 \\
0.10 & 0.533 & 0.543 & 0.538 & +1.0 & +0.5 \\
0.20 & 0.524 & 0.531 & 0.527 & +0.7 & +0.3 \\
\bottomrule
\end{tabular}
\caption{Temperature ablation on SWE-smith validation (Qwen, R@1). Embedding benefit direction is consistent across all $\tau$ values.}
\label{tab:temp_qwen}
\end{table}

\begin{table}[t]
\centering
\small
\begin{tabular}{lccccc}
\toprule
$\tau$ & \textbf{Random} & \textbf{Embed} & \textbf{Dataflow} & \textbf{Embed $\Delta$} & \textbf{Dataflow $\Delta$} \\
\midrule
0.05 & 0.438 & 0.443 & 0.430 & +0.5 & $-$0.8 \\
0.10 & 0.449 & 0.456 & 0.441 & +0.7 & $-$0.8 \\
0.20 & 0.442 & 0.446 & 0.435 & +0.4 & $-$0.7 \\
\bottomrule
\end{tabular}
\caption{Temperature ablation on SWE-smith validation (StarCoder, R@1). Dataflow consistently underperforms random on this encoder.}
\label{tab:temp_starcoder}
\end{table}

The embedding $>$ random direction is robust across all temperature settings for both encoders, supporting the claim that embedding-distance mining provides genuine benefit independent of the temperature regime. Dataflow shows inconsistent direction across encoders: positive for Qwen at all temperatures but negative for StarCoder at all temperatures. The optimal temperature is $\tau = 0.1$ for both encoders, with $\tau = 0.05$ increasing cross-seed variance and $\tau = 0.2$ flattening the difficulty gradient.

%------------------------------------------------------------
\section{Ablation Studies}
\label{sec:ablations}
%------------------------------------------------------------

\subsection{Call-Graph Hop Depth}
\label{sec:hops}

We ablate the hop depth for dataflow connectivity (Table~\ref{tab:hops}). The 2-hop setting slightly outperforms 1-hop (+0.7pp) as it captures transitive dependencies, but 3-hop degrades ($-$0.4pp vs.\ 2-hop) because deep connectivity in large repositories approaches full-graph reachability, diluting the structural signal. Coverage also decreases with depth as fewer disconnected negatives remain.

\begin{table}[t]
\centering
\small
\begin{tabular}{lccc}
\toprule
\textbf{Hops} & \textbf{R@1} & \textbf{Coverage (\%)} & \textbf{Avg neg/query} \\
\midrule
1-hop & 0.531 & 42.1 & 8.3 \\
2-hop & 0.538 & 34.3 & 5.7 \\
3-hop & 0.534 & 28.7 & 3.9 \\
\bottomrule
\end{tabular}
\caption{Hop depth ablation (Qwen, dataflow strategy). Differences are within noise range (cross-seed std $\approx$ 1.5--3.0pp).}
\label{tab:hops}
\end{table}

The differences are within the cross-seed variance ($\sim$1.5--3.0pp), so we treat 2-hop as a reasonable default rather than a critical design choice.

\subsection{Training Data Size Scaling}
\label{sec:scaling}

Table~\ref{tab:scaling} reports R@1 as a function of training set size. Overall performance improves from 1K to 3K instances but plateaus at 5K, consistent with overfitting observed after step 1000 in training runs.

\begin{table}[t]
\centering
\small
\begin{tabular}{lcccc}
\toprule
\textbf{Train Size} & \textbf{Random R@1} & \textbf{Embed R@1} & \textbf{$\Delta$ (pp)} \\
\midrule
1K & 0.500 & 0.518 & +1.8 \\
3K & 0.530 & 0.541 & +1.1 \\
5K & 0.529 & 0.538 & +0.9 \\
\bottomrule
\end{tabular}
\caption{Data scaling for Qwen (embedding strategy). The embedding benefit diminishes with more training data.}
\label{tab:scaling}
\end{table}

The embedding benefit diminishes from +1.8pp at 1K to +0.9pp at 5K, suggesting that larger datasets provide sufficient coverage to reduce dependence on mining strategy quality. This is a relevant limitation: the practical value of sophisticated mining may diminish at production data scales.

\subsection{Fusion Head}
\label{sec:fusion}

Table~\ref{tab:fusion} compares three scoring heads: embedding-only, structural-features-only, and combined.

\begin{table}[t]
\centering
\small
\begin{tabular}{lcc}
\toprule
\textbf{Head} & \textbf{Qwen R@1} & \textbf{StarCoder R@1} \\
\midrule
Embed-only & 0.538 & 0.451 \\
Struct-only & 0.489 & 0.412 \\
Combined & 0.541 & 0.448 \\
\bottomrule
\end{tabular}
\caption{Fusion head ablation. Structural features alone are weak; the combined head provides marginal additive benefit for Qwen (+0.3pp) but slightly hurts StarCoder ($-$0.3pp).}
\label{tab:fusion}
\end{table}

Structural features alone perform poorly (R@1 = 0.489/0.412) because CodeQL dataflow features are binary and sparse, lacking fine-grained discrimination. The combined head adds only +0.3pp over embed-only for Qwen and is slightly harmful for StarCoder ($-$0.3pp), suggesting that the structural signal's primary value is in \emph{negative selection} during training, not in scoring at inference time.

\subsection{Weighted vs.\ Binary Dataflow Negatives}
\label{sec:weighted}

We test whether distance-weighted sampling from the dataflow graph improves over binary (connected/disconnected) selection (Table~\ref{tab:weighted}).

\begin{table}[t]
\centering
\small
\begin{tabular}{lcc}
\toprule
\textbf{Method} & \textbf{Qwen R@1} & \textbf{StarCoder R@1} \\
\midrule
Binary (default) & 0.538 & 0.448 \\
Distance-weighted & 0.535 & 0.444 \\
\bottomrule
\end{tabular}
\caption{Binary vs.\ distance-weighted dataflow negative selection. No improvement from weighting.}
\label{tab:weighted}
\end{table}

Distance-weighted sampling does not improve over binary selection, likely due to the bimodal distribution of dataflow distances (most pairs are either 1--2 hops or completely disconnected), which causes smooth weighting to degenerate to approximate binary selection.

\subsection{Static vs.\ Dynamic Mining}
\label{sec:dynamic}

Table~\ref{tab:dynamic} compares static mining (once from pretrained encoder) with dynamic re-mining per epoch.

\begin{table}[t]
\centering
\small
\begin{tabular}{lccc}
\toprule
\textbf{Mining} & \textbf{Qwen R@1} & \textbf{StarCoder R@1} & \textbf{GPU hours} \\
\midrule
Static & 0.543 & 0.456 & 3.5 \\
Dynamic & 0.547 & 0.452 & 11.2 \\
\bottomrule
\end{tabular}
\caption{Static vs.\ dynamic mining. Dynamic re-mining provides no significant benefit at $3.2\times$ compute cost.}
\label{tab:dynamic}
\end{table}

Dynamic re-mining does not significantly improve over static mining ($\Delta = +0.4$pp Qwen, $-$0.4pp StarCoder), which is expected given the limited representation drift in 3-epoch LoRA fine-tuning. The $3.2\times$ compute overhead is not justified. This also confirms that the static mining protocol used throughout our experiments is not a confound.

%------------------------------------------------------------
\section{Qualitative Analysis}
\label{sec:qualitative}
%------------------------------------------------------------

We curate 8 representative cases where structural negatives changed retrieval outcomes (4 positive, 2 negative, 2 ambiguous).

\paragraph{Positive examples.} In 4 cases, dataflow-disconnected negatives help the retriever distinguish lexically similar but structurally unrelated files. For instance, in a Django repository, two utility modules share significant token overlap (\texttt{utils/validators.py} and \texttt{utils/formatters.py}) but have no dataflow connection. Dataflow-aware training teaches the encoder to separate these, improving retrieval of the true bug-fix target.

\paragraph{Negative examples.} In 2 cases, the structural bias introduced by dataflow negatives causes misses. One involves a file pair connected through a rarely-exercised code path that CodeQL fails to extract; the model learns to treat them as disconnected when they are functionally related.

\paragraph{Ambiguous examples.} Two cases show mixed signals where the retrieval outcome differs between seeds, suggesting the structural signal interacts with other training noise.

The pattern across cases is not clean enough to support strong causal inference from $n = 2$ encoders, but the positive examples demonstrate a plausible mechanism by which structural negatives can complement text-based mining.

%------------------------------------------------------------
\section{Discussion}
\label{sec:discussion}
%------------------------------------------------------------

\subsection{Positioning Relative to Inference-Time Methods}

Table~\ref{tab:inference_compare} contextualizes our approach against inference-time structural methods.

\begin{table}[t]
\centering
\small
\begin{tabular}{lllll}
\toprule
\textbf{Method} & \textbf{Type} & \textbf{Structural Signal} & \textbf{Inf.\ Cost} & \textbf{Reported Gain} \\
\midrule
StructFuse (ours) & training & CodeQL dataflow & zero & +0.5--1.5pp R@1 \\
CodexGraph & inference & LLM-built graph & high & +5.2pp R@1 \\
LocAgent & inference & call graph nav. & moderate & 26.3\% resolve \\
\bottomrule
\end{tabular}
\caption{Training-time vs.\ inference-time structural methods. StructFuse is complementary: it improves the base retriever at zero inference cost.}
\label{tab:inference_compare}
\end{table}

StructFuse operates at training time with zero inference overhead, improving the base retriever that inference-time methods build upon. The gains are modest compared to inference-time approaches, but the two paradigms are complementary: a better base retriever provides a stronger starting point for graph-guided navigation. Validating this complementarity is an important direction for future work.

\subsection{Limitations}

\begin{itemize}
\item \textbf{Two encoders.} Our encoder-strategy interaction observation is based on $n = 2$ architectures. While the representation geometry mechanism provides explanatory power, broader validation is needed.
\item \textbf{Diminishing returns at scale.} The embedding mining benefit shrinks with more training data (Table~\ref{tab:scaling}), suggesting limited practical value at production scales.
\item \textbf{Downstream gap.} Retrieval-level improvements do not translate to statistically significant downstream gains on SWE-bench Lite. The generation bottleneck dominates.
\item \textbf{CodeQL coverage.} Dataflow extraction covers 123/131 repos but structural relationships are incomplete---CodeQL misses dynamically dispatched calls and reflection-heavy code.
\item \textbf{Confidence intervals.} For RepoBench-R, the StarCoder embedding benefit (+0.5pp) has a 95\% CI that includes zero. We characterize all claims as directional evidence, not definitive conclusions.
\end{itemize}

%------------------------------------------------------------
\section{Conclusion}
\label{sec:conclusion}
%------------------------------------------------------------

We present a controlled study of hard negative mining strategies for contrastive code retrieval. Our findings provide directional evidence that true embedding-distance mining tends to improve retrieval over random baselines, correcting prior reports based on BM25-proxy confounds. Dataflow-grounded negatives show conditional, encoder-dependent effects that we link to representation geometry---an exploratory observation warranting further investigation.

Our primary contributions are methodological: a controlled comparison framework that isolates mining strategy effects, a representation geometry analysis that explains encoder-strategy interactions, and a scalable CodeQL extraction pipeline for structural negatives. End-to-end evaluation on SWE-bench Lite yields no significant differences, an important finding for calibrating expectations about retrieval-only improvements in agent pipelines.

These results suggest that the next frontier lies not in harder negatives alone, but in generation-aware retrieval training that accounts for downstream bottlenecks. We release our framework and data to support this investigation.

\bibliography{custom}

\end{document}
